\documentclass[10pt,a4paper]{amsart}
\usepackage[margin=19mm]{geometry}
\usepackage{amsmath,amssymb,mathtools}
\usepackage[scr=rsfs,cal=euler]{mathalfa}
\usepackage{xfrac}
\usepackage[unicode,hidelinks,bookmarks=false]{hyperref}
\newcommand{\ie}{i.e.\ }
\newcommand{\cf}{cf.\ }
\newcommand{\resp}{respectively\ }
\newcommand{\Subsection}{\subsection}
\providecommand{\nobreakdash}{\nobreak\hbox{-}}
\setlength{\emergencystretch}{2em}
\multlinegap0pt
\usepackage{bm}
\usepackage{bbm}
\usepackage{dsfont}
\usepackage{xspace}
\usepackage{cancel}
\usepackage{mathtools}
\usepackage{amsthm}
\makeatletter
\@ifundefined{theorem}{\newtheorem{theorem}{Theorem}}{}
\makeatother
\makeatletter
\expandafter\ifx\csname example\endcsname\relax
\newenvironment{example}
{\pushQED{\qed}\renewcommand{\qedsymbol}{$\diamond$}\examplex}
{\popQED\endexamplex}
\fi
\makeatother
\title[Positive Polytopes with Few Facets in the Grassmannian]{Positive Polytopes with Few Facets in the Grassmannian}
\author{Dmitrii Pavlov, Kristian Ranestad}
\date{Source: arXiv:2503.01652v2 (CC BY 4.0)}
\begin{document}
\maketitle

\noindent\emph{Source and licence.} Positive Polytopes with Few Facets in the Grassmannian. Version arXiv:2503.01652v2 (CC BY 4.0), distributed under the Creative Commons Attribution 4.0 International licence, as recorded in the arXiv licence metadata for this exact version and shown on the abstract page https://arxiv.org/abs/2503.01652v2. Full rights notice: licenses/arxiv-2503.01652-CC-BY-4.0.txt.\par\smallskip

\noindent\textit{Editorial scope.} Counted unit: the Theorem whose source label is \texttt{thm:penta} in the article (source0.tex lines 258--260, source label \texttt{thm:penta}), delivered together with the article's own complete original proof of it (source0.tex lines 262--284, one \texttt{proof} environment of 23 lines). No mathematics is written, repaired or re-derived: statement and proof are the article's own text line for line. Required context retained uncounted: none. Every definition, display and label of those lines is retained unchanged. Omitted from this package and not redistributed: 1--257, 261--261, 285--711. Nothing inside a retained excerpt is omitted or altered: every retained body is delivered exactly as the article's lines are written, so no omission receipt of any kind accompanies this package. Citations: the retained excerpts carry no citation command at all, so no bibliography entry of the article is transcribed and no reference is redistributed. Reference adaptation: none was needed and none was made; every reference inside the delivered unit resolves inside it, so no unresolved reference or citation is produced. Printed slips kept literal: no printed word, symbol, hypothesis or number of the counted statement or of its proof was altered. Layout and class: the article's own class is \texttt{extarticle} with packages including amsmath, amsthm, amssymb, color, hyperref, graphicx, caption, mathtools, enumerate, verbatim, tikz, tikz-cd, tikz-3dplot, algorithm; the wrapper is the lane's pinned \texttt{amsart} editorial wrapper at 10pt on A4, which restates the theorem-like environments the retained text needs and adds the article's own macro definitions that the retained text needs (none), with extra wrapper packages (bm, bbm, dsfont, xspace, cancel, mathtools). The delivered numbering is the unit's own. Assets: no retained span contains a figure environment, a graphics-inclusion command, a PDF-page inclusion, a table or a TikZ picture; no image file is copied into this package, the package declares no asset dependency and no image file is redistributed with it. Licence: version arXiv:2503.01652v2 of this article is distributed under the Creative Commons Attribution 4.0 International licence, as recorded in the arXiv licence metadata for this exact version and shown on the abstract page https://arxiv.org/abs/2503.01652v2. A full rights notice is delivered at licenses/arxiv-2503.01652-CC-BY-4.0.txt. Characters: every retained excerpt is pure ASCII, so the delivered engine, XeLaTeX with its default Latin Modern text font, reports no missing character.
\begin{theorem} \label{thm:penta}
    Suppose $c_I\neq 0$ for all $I$, and suppose the hyperplane $H$ defined by $h(\mathbf{p}) = 0$ intersects every facet of $\mathrm{Gr}_{\geq 0}(2,4)$. Then there exists a unique adjoint of the positive pentahedron $S = P \cap \mathrm{Gr}(2,4)$, and $S$ is a positive geometry. The numerator of its canonical form (and the defining polynomial of the adjoint) is $\sum \overline{c_I}p_I$, where $\overline{c_I} = \max (c_I, 0)$.
\end{theorem}

\begin{proof}
    Let $S_1$ and $S_2$ be the pieces into which $\mathrm{Gr}_{\geq 0}(2,4)$ is cut by $H$. In the notation above, we have $S_1 = S$ and $S_2 = (\mathrm{Gr}_{\geq 0}(2,4) \setminus S) \cup (\mathrm{Gr}_{\geq 0}(2,4) \cap H)$. We will now describe the residual arrangements $\mathcal{R}_1$ and $\mathcal{R}_2$ of $S_1$ and $S_2$, respectively. By the first assumption, each facet of $\mathrm{Gr}_{\geq 0}(2,4)$ contributes a facet to both $S_1$ and $S_2$. The residual arrangement $\mathcal{R}_i$ consists of two parts. The first one (we denote it $\mathcal{R}^G_i$) consists of the intersections of facet hyperplanes of $\mathrm{Gr}_{\geq 0}(2,4)$ that do not give a face of $S_i$. The second one (we denote in $\mathcal{R}^H_i$) consists of the intersections of $H$ with some of the facet hyperplanes of $\mathrm{Gr}_{\geq 0}(2,4)$ that do not give a face of $S_i$. Since $c_I\neq 0$ for all $I$ by assumption, $H$ does not contain any vertex of $\mathrm{Gr}_{\geq 0}(2,4)$. This means that each vertex of $\mathrm{Gr}_{\geq 0}(2,4)$ is either in $\mathcal{R}^G_{1}$ or in $\mathcal{R}^G_2$. More precisely, $v_I \in \mathcal{R}^G_{1}$ if $c_I <0$ and $v_I \in \mathcal{R}^G_2$ if $c_I > 0$. Note that if for two vertices $v$ and $w$ connected by an edge of $\mathrm{Gr}_{\geq 0}(2,4)$ we have $v \in \mathcal{R}^G_{1}$ and $w \in \mathcal{R}^G_2$, then the edge of $\mathrm{Gr}_{\geq 0}(2,4)$ between them contributes an edge to both $S_1$ and $S_2$, and is in none of the residual arrangements. This means that the span of $\mathcal{R}^G_i$ is equal to the span of the vertices of $\mathrm{Gr}_{\geq 0}(k,n)$ that are in $\mathcal{R}^G_i$. 

We now turn to characterizing $\mathcal{R}^H_i$. If $F$ is a face of $\mathrm{Gr}_{\geq 0}(2,4)$ that contributes a face to both $S_1$ and $S_2$, then it has to intersect $H$ inside $\mathrm{Gr}_{\geq 0}(2,4)$, and therefore $F\cap H$ does not contribute to either of the residual arrangements. If $F$ only contributes a face to $S_1$, then it is in $\mathcal{R}^G_2$ and the hyperplane spanned by $F$ intersects $H$ outside of $\mathrm{Gr}_{\geq 0}(2,4)$. This intersection is therefore in $\mathcal{R}^H_1$. If $F$ only contributes a face to $S_2$, then it already is in $\mathcal{R}^G_1$ and only contributes a lower-dimensional piece to $\mathcal{R}^H_1$. Since each face $F$ of $\mathrm{Gr}_{\geq 0}(2,4)$ contributes a face to etiher $S_1$ or $S_2$, and since all the irreducible components of $\mathcal{R}_i^H$ are by definition Zariski closures of the sets of the form $F \cap H$ for some face $F$ of $\mathrm{Gr}_{\geq 0}(2,4)$, one sees that $\mathrm{span}(\mathcal{R}^H_1) = \mathrm{span}(H\cap \mathcal{R}^G_2)$ and, symmetrically, $\mathrm{span}(\mathcal{R}^H_2) = \mathrm{span}(H\cap \mathcal{R}^G_1)$. Since $\mathcal{R}_i = \mathcal{R}_i^G \cup \mathcal{R}_i^H$, we now can describe the span of $\mathcal{R}_i$, which is a projective space.

Since the span of $\mathcal{R}_2^G$ is equal to the span of its vertices, the span of $H \cap \mathcal{R}_2^G$ is equal to the span of the intersection points $w_j$ of the edges between these vertices with $H$.
Any hyperplane interpolating $\mathcal{R}_1^G$ has an equation of the form 
$$\sum\limits_{I \in \binom{[n]}{k}} b_I p_I,$$
where $b_I =0$ if $v_I \in \mathcal{R}_1$. The fact that the adjoint hyperplane should also interpolate the points $w_j$ imposes the conditions $b_I = Ac_I$ if $v_I\in S_1$ for the adjoint, where $A$ is some scalar factor. This is because each point $w_j$ is of the form $\alpha v_I + \beta v_J$ for the vertices $v_I$ and $v_J$ in $\mathcal{R}_G^2$. In addition, for any hyperplane interpolating $w_j$ we have $0=\alpha b_I+ \beta b_J$, and, since $H$ contains $w_j$, we also have $0=\alpha c_I+ \beta c_J$ for all $I$ and $J$ such that $v_I, v_J \in \mathcal{R}_G^2$. This determines the equation of the adjoint $a(\mathbf{p})$ uniquely up to a scalar factor. This scalar factor is then fixed by computing the residue of the resulting canonical form at any vertex of $S_1$ and requiring it to be $\pm 1$ depending on the chosen orientation.

The conditions $b_I=c_I$ for $v_I\in S_1$ can alternatively be obtained by imposing the residue conditions for the canonical form at the vertices $v_I \in S_1$. This shows that the form 
$$
\Omega(\mathrm{Gr}(2,4), S_1) = \dfrac{a(\mathbf{p}) dp_{12}\wedge dp_{23} \wedge dp_{34} \wedge dp_{14} }{p_{12}p_{23}p_{34}p_{14} h(\mathbf{p})}.
$$
is the canonical form of the positive pentahedron $S_1$ and that $S_1$ is a positive geometry. Analogously one can show that $S_2$ is also a positive geometry.
% The span of $H \cap \mathcal{R}^G_2$ is contained in $H\cap \mathrm{span}(\{v_I:c_I<0\})$. The ideal of this latter intersection contains the defining polynomial $\sum c_I p_I$ of $H$ as well as any linear form 
%$$\sum\limits_{I \in \binom{[n]}{k}} b_I p_I,$$
%where $b_I = 0$ if $v_I \in S_1$. In particular, by setting $b_I = c_I$ for $v_I\in \mathcal{R}_1$ and subtracting the resulting form from $\sum c_I p_I$, we get the defining equation of $A$. Thus, $A$ interpolates $\mathcal{R}^H_1$. It therefore interpolates the whole residual arrangement $\mathcal{R}_1$ and is the adjoint of $S_1$. The canonical form of $S_1$ is then given by 
%$$
%\Omega(\mathrm{Gr}(2,4), S_1) = \dfrac{A}{p_{12}p_{23}p_{34}p_{14} h(\mathbf{p})}.
%$$

\end{proof}

\end{document}
